\textbf{Phần 1. Trắc nghiệm lựa chọn}
\setcounterpageref{ex}{0}

\begin{ex} %[2D8N39]
Kết quả của phép tính $2 \times 4$ là
\choice
{$6$}
{$10$}
{$12$}
{\True $8$}
\loigiai{
$2 \times 4 = 8$.\\
\textbf{Đáp án của câu hỏi là D}
}
\end{ex}

\begin{ex} %[2D8N39]
Kết quả của phép tính $2 \times 7$ là
\choice
{$12$}
{$16$}
{\True $14$}
{$18$}
\loigiai{
$2 \times 7 = 14$.\\
\textbf{Đáp án của câu hỏi là C}
}
\end{ex}

\begin{ex} %[2D8H39]
Phép tính nào sau đây có kết quả bằng 16?
\choice
{\True $2 \times 8$}
{$2 \times 7$}
{$2 \times 9$}
{$2 \times 6$}
\loigiai{
$2 \times 8 = 16$.\\
\textbf{Đáp án của câu hỏi là A}
}
\end{ex}

\begin{ex} %[2D8H39]
Phép tính nào trong bảng nhân 2 có kết quả lớn nhất?
\choice
{$2 \times 8$}
{$2 \times 9$}
{$2 \times 7$}
{\True $2 \times 10$}
\loigiai{
$2 \times 10 = 20$ (lớn nhất).\\
\textbf{Đáp án của câu hỏi là D}
}
\end{ex}

\begin{ex} %[2D8H39]
Trong bảng nhân 2, hai tích liên tiếp hơn kém nhau bao nhiêu đơn vị?
\choice
{\True $2$ đơn vị}
{$1$ đơn vị}
{$3$ đơn vị}
{$4$ đơn vị}
\loigiai{
Hai tích liên tiếp hơn kém nhau 2 đơn vị.\\
\textbf{Đáp án của câu hỏi là A}
}
\end{ex}

\begin{ex} %[2D8H39]
Điền số thích hợp: $2 \times \dots = 18$.
\choice
{$8$}
{$7$}
{\True $9$}
{$6$}
\loigiai{
$2 \times 9 = 18$.\\
\textbf{Đáp án của câu hỏi là C}
}
\end{ex}

\begin{ex} %[2D8V39]
Mỗi bạn học sinh có 2 hộp bút. Hỏi 8 bạn học sinh có tất cả bao nhiêu hộp bút?
\choice
{$14$ hộp}
{\True $16$ hộp}
{$18$ hộp}
{$12$ hộp}
\loigiai{
$2 \times 8 = 16$ hộp.\\
\textbf{Đáp án của câu hỏi là B}
}
\end{ex}

\begin{ex} %[2D8V39]
Mỗi đôi đũa có 2 chiếc đũa. Hỏi 9 đôi đũa có tất cả bao nhiêu chiếc đũa?
\choice
{$16$ chiếc}
{$20$ chiếc}
{$14$ chiếc}
{\True $18$ chiếc}
\loigiai{
$2 \times 9 = 18$ chiếc đũa.\\
\textbf{Đáp án của câu hỏi là D}
}
\end{ex}

\begin{ex} %[2D8V39]
Tính giá trị biểu thức: $2 \times 6 + 18$.
\choice
{$28$}
{\True $30$}
{$32$}
{$34$}
\loigiai{
$12 + 18 = 30$.\\
\textbf{Đáp án của câu hỏi là B}
}
\end{ex}

\begin{ex} %[2D8V39]
Điền dấu thích hợp: $2 \times 5 + 6 \dots 2 \times 8$.
\choice
{$>$}
{$<$}
{\True $=$}
{Không so sánh được}
\loigiai{
$2 \times 5 + 6 = 10 + 6 = 16$ và $2 \times 8 = 16$. Vậy bằng nhau.\\
\textbf{Đáp án của câu hỏi là C}
}
\end{ex}

\begin{ex} %[2D8V39]
Có 7 can nước mắm, mỗi can chứa $2\text{ l}$. Cả 7 can chứa số lít nước mắm là
\choice
{$12\text{ l}$}
{\True $14\text{ l}$}
{$16\text{ l}$}
{$18\text{ l}$}
\loigiai{
$2 \times 7 = 14\text{ l}$.\\
\textbf{Đáp án của câu hỏi là B}
}
\end{ex}

\begin{ex} %[2D8C39]
Tìm số thích hợp điền vào ô trống: $2 \times \dots > 15$ và $< 18$.
\choice
{\True $8$}
{$7$}
{$9$}
{$6$}
\loigiai{
$2 \times 8 = 16$ thỏa mãn $15 < 16 < 18$.\\
\textbf{Đáp án của câu hỏi là A}
}
\end{ex}

\textbf{Phần 2. Trắc nghiệm đúng sai}
\setcounterpageref{ex}{0}

\begin{ex} %[2D8H39]
Xét tính đúng sai của các phép tính trong bảng nhân 2:
\choiceTF
{\True $2 \times 3 = 6$.}
{\True $2 \times 7 = 14$.}
{$2 \times 9 = 19$.}
{\True $2 \times 9 = 18$.}
\loigiai{
\begin{itemize}
    \item $2 \times 3 = 6$: Đúng.
    \item $2 \times 7 = 14$: Đúng.
    \item $2 \times 9 = 18 \neq 19$. Ý 3 Sai, Ý 4 Đúng.
\end{itemize}
\begin{itemize}
    \item a) \textbf{Đúng}.
    \item b) \textbf{Đúng}.
    \item c) \textbf{Sai}.
    \item d) \textbf{Đúng}.
\end{itemize}
}
\end{ex}

\begin{ex} %[2D8V39]
Một đôi chim bồ câu có 2 con chim. Xét tính đúng sai:
\choiceTF
{\True 6 đôi chim có 12 con chim ($2 \times 6 = 12$).}
{\True 10 đôi chim có 20 con chim ($2 \times 10 = 20$).}
{8 đôi chim có 14 con chim.}
{\True 8 đôi chim có $2 \times 8 = 16$ con chim.}
\loigiai{
\begin{itemize}
    \item 6 đôi: $2 \times 6 = 12$: Đúng.
    \item 10 đôi: $2 \times 10 = 20$: Đúng.
    \item 8 đôi: $2 \times 8 = 16 \neq 14$. Ý 3 Sai, Ý 4 Đúng.
\end{itemize}
\begin{itemize}
    \item a) \textbf{Đúng}.
    \item b) \textbf{Đúng}.
    \item c) \textbf{Sai}.
    \item d) \textbf{Đúng}.
\end{itemize}
}
\end{ex}

\textbf{Phần 3. Trả lời ngắn}
\setcounterpageref{ex}{0}

\begin{ex} %[2D8H39]
Tính: $2 \times 5$.
\shortans{10}
\loigiai{
\textbf{Đáp số:} $10$\\
$2 \times 5 = 10$.
}
\end{ex}

\begin{ex} %[2D8H39]
Tính: $2 \times 8$.
\shortans{16}
\loigiai{
\textbf{Đáp số:} $16$\\
$2 \times 8 = 16$.
}
\end{ex}

\begin{ex} %[2D8V39]
Tính giá trị biểu thức: $2 \times 9 - 8$.
\shortans{10}
\loigiai{
\textbf{Đáp số:} $10$\\
$18 - 8 = 10$.
}
\end{ex}

\begin{ex} %[2D8V39]
Mỗi con vịt có 2 cái chân. Hỏi 10 con vịt có bao nhiêu cái chân?
\shortans{20}
\loigiai{
\textbf{Đáp số:} $20$\\
$2 \times 10 = 20$ cái chân.
}
\end{ex}

\begin{ex} %[2D8V39]
Điền số thích hợp: $2 \times ? = 14$.
\shortans{7}
\loigiai{
\textbf{Đáp số:} $7$\\
$14 : 2 = 7$.
}
\end{ex}

\begin{ex} %[2D8C39]
Có 6 túi kẹo, mỗi túi có 2 cái kẹo. Mẹ mua thêm 8 cái kẹo nữa. Hỏi có tất cả bao nhiêu cái kẹo?
\shortans{20}
\loigiai{
\textbf{Đáp số:} $20$\\
$2 \times 6 + 8 = 12 + 8 = 20$ cái.
}
\end{ex}
